% P16-C1: finite propositional abduction and explicit input boundaries.
\chapter{Finite Logical Abduction}
\label{sec:p16-finite-abduction}

Logical abduction searches for permitted assumptions which, together with a
background theory, consistently entail specified observations. This chapter
uses Eiter and Gottlob's finite propositional model: hypotheses and
observations are positive atoms, and an explanation is a subset of the
permitted hypotheses \cite[sec.~2, manuscript pp.~9--11]{EiterGottlobAbductionManuscript}.
The finite construction and examples below are supplied here with complete
proofs. No complexity classification is needed for this construction.

The source requires a consistent background theory. We also allow an
inconsistent theory as an input to the checking procedure, and prove that
this additional input case has no explanations. This extension does not
change the explanation criterion on the source's inputs.

\section{Background, hypotheses and observations}
\label{sec:p16-abduction-inputs}

\begin{hbdefinition}{def:p16-abduction-input}{Finite abduction input}
A \emph{finite abduction input} is a tuple
\(\mathcal P=(V,H,M,T)\), with the following data:
\begin{enumerate}
\item \(V\) is a finite subset of the propositional alphabet \(\mathsf{Prop}\).
\item \(H\subseteq V\) is the set of \emph{permitted hypotheses}, also
      called the \emph{abducible atoms}.
\item \(M\subseteq V\) is the set of \emph{observations}; Eiter and
      Gottlob call these \emph{manifestations}.
\item \(T\) is a finite set of classical propositional formulas whose
      occurring atoms all belong to \(V\), called the \emph{background
      theory}.
\end{enumerate}
The sets may be empty. The name \emph{abducible} is borrowed: Eiter and
Gottlob remark that a hypothesis of their model does the work of what
other models call an abducible proposition, and refer for those models to
Console, Theseider Dupr\'e and Torasso's paper on the relationship between
abduction and deduction (their reference~14, Journal of Logic and
Computation, 1991)
\cite[after def.~2.2, manuscript p.~10, and reference list, manuscript p.~41]{EiterGottlobAbductionManuscript}.
Neither \(H\cap M=\varnothing\) nor
\(H\cup M=V\) is required. Eiter and Gottlob give a reason for each
omission. Overlap is allowed because one observation may explain another.
Hypotheses and observations need not exhaust \(V\) because explanations
are assumption-based: only some propositions may occur in an explanation,
which suits a search that focuses attention on a chosen subset of facts.
Their further remark concerns the scope of their computational results,
not the explanation relation: omitting the two restrictions leaves no gap
in those results, which remain valid when either restriction is imposed.
The grounds they give are about their own proofs. The condition
\(H\cap M=\varnothing\) holds in all but a few of those proofs, the
condition \(H\cup M=V\) is assumed in all of them with rare exceptions,
and they state that the remaining proofs are easily adapted
\cite[after def.~2.2, manuscript p.~10]{EiterGottlobAbductionManuscript}.
This chapter does not reproduce those computational results, so the remark
is reported here as the source's claim about its proofs.
Hypotheses and observations in this definition
are atoms themselves: negative literals and compound formulas are not
members of \(H\) or \(M\). For a similar model, the source points to
Konolige's causal theories (its reference~[40], Konolige's 1992 paper
on abduction and closure in causal theories). It states that the restriction
to positive atoms loses no
generality, since a formula \(F\) becomes available as a hypothesis when
a new variable \(x'\) is added to the hypotheses together with the
formula \(x'\leftrightarrow F\) in the theory
\cite[manuscript p.~10]{EiterGottlobAbductionManuscript}.
This construction keeps every atom of \(H\) as a hypothesis and adds
\(x'\) beside them. \Cref{thm:p15-w02c-p15-027-source-extension} proves
the exact sense in which it loses no generality: the explanations of the
extended input are the old explanations together with the sets
\(S\cup\{x'\}\) for which \(S\cup\{F\}\) explains \(M\), while
minimum-cardinality explanations can change.
\Cref{thm:p15-w02c-p15-027-formula-hypotheses} treats the different
construction in which every permitted hypothesis is a formula and each is
replaced by a new atom.
No special syntactic form of \(T\) is required. The source adds that
implementations often put the theory into a special format, most commonly
clausal form, and that clausal form appears not to affect the
computational complexity of logic-based abduction, which it calls
unsurprising \cite[manuscript p.~10]{EiterGottlobAbductionManuscript};
its summary of results says that clausal theories have exactly the same
complexity results as general theories
\cite[manuscript p.~8]{EiterGottlobAbductionManuscript}. That complexity
claim is the source's. The semantic fact used in this chapter, that
replacing \(T\) by an equivalent clausal theory over \(V\) leaves every
explanation family unchanged, is
\cref{thm:p15-w02c-p15-027-theory-forms}.
The theory \(T\) may contain negations
and compound formulas. Its formulas use the grammar and abbreviations of
\cref{sec:prop-syntax-semantics-01}.
The source's truth constant and biconditional are written here using
the book's abbreviations \(\topv\) and \(\leftrightarrow\).

A \emph{candidate} is a set \(E\subseteq H\). In \(T\cup E\),
the atoms in \(E\) are treated as formulas asserted positively.
An atom outside \(H\) cannot be inserted into a candidate merely because
it would imply an observation. Changing \(H\) changes the input problem.
\index{abduction!permitted hypotheses}\index{abduction!observations}
\end{hbdefinition}

\begin{hbdefinition}{def:p16-abduction-models}{Finite valuations and explanatory entailment}
Use the truth values \(\mathbb{B}=\{\mathbf{F},\mathbf{T}\}\) and the recursive
truth clauses of \cref{sec:prop-syntax-semantics-02}. Put
\[
 \mathsf{Val}(V)\defeq\{v:v\text{ is a function }V\to\mathbb{B}\}.
\]
For \(v\in\mathsf{Val}(V)\), evaluate a formula over \(V\) by those
same clauses, and write \(v\entails\varphi\) when its value is
\(\mathbf{T}\). For a set \(\Gamma\) of formulas over \(V\), put
\[
 \operatorname{Mod}_V(\Gamma)
   \defeq\{v\in\mathsf{Val}(V):
                  v\entails\gamma\text{ for every }\gamma\in\Gamma\}.
\]
We say \(\Gamma\) is \emph{consistent} in this chapter when this set
is nonempty, that is, when \(\Gamma\) is satisfiable. The notation
\(\Gamma\entails M\) abbreviates the assertion
\(\Gamma\entails m\) for every \(m\in M\), using the book's
classical semantic consequence relation. The next theorem identifies this
relation with the finite valuation test.

A candidate \(E\) is an \emph{explanation} for \(\mathcal P\) when
\[
 \operatorname{Mod}_V(T\cup E)\ne\varnothing
       \qquad\text{and}\qquad T\cup E\entails M.
\]
Write \(\operatorname{Expl}(\mathcal P)\) for the family of all such
explanations. This family corresponds to the source's
\(\operatorname{Sol}(\mathcal P)\).
\index{abduction!explanation}\index{abduction!consistency}
\end{hbdefinition}

\begin{hbtheorem}{thm:p16-abduction-semantics}{proved=annex:p16.abduction-semantics}{The finite criterion and its boundary cases}
For every finite abduction input and every \(E\subseteq H\),
\(E\) is an explanation exactly when
\(\operatorname{Mod}_V(T\cup E)\ne\varnothing\) and
\[
 v(m)=\mathbf{T}\qquad
  (v\in\operatorname{Mod}_V(T\cup E),\ m\in M).
\]
The use of functions on \(V\) agrees with the book's valuations on all
of \(\mathsf{Prop}\). After naming the finite alphabet in \(\mathsf{Prop}\) and
expanding those two abbreviations, inputs with consistent \(T\) have
precisely Eiter and Gottlob's explanation criterion
\cite[defs.~2.1--2.2, manuscript pp.~9--10]{EiterGottlobAbductionManuscript}.
The additional and empty cases satisfy:
\begin{enumerate}
\item If \(T\) is inconsistent, then
      \(\operatorname{Expl}(\mathcal P)=\varnothing\).
\item The empty set is an explanation exactly when \(T\) is consistent
      and \(T\entails M\).
\item If \(M=\varnothing\), the explanations are exactly those
      \(E\subseteq H\) for which \(T\cup E\) is consistent.
\end{enumerate}
\end{hbtheorem}

\begin{hbdefinition}{def:p15-w02c-p15-027-theory-forms}{Forms of a background theory}
Eiter and Gottlob fix the following vocabulary for theories
\cite[manuscript p.~9]{EiterGottlobAbductionManuscript}. The clauses are
the book's clauses of \cref{sec:prop-syntax-semantics-04}, finite sets of
literals read as disjunctions.
\begin{enumerate}
\item A \emph{literal} is an atom or a negated atom. A \emph{clause} is a
      disjunction
      \(x_1\lor\cdots\lor x_k\lor\neg x_{k+1}\lor\cdots\lor\neg x_n\)
      of pairwise distinct literals. Clauses that differ only in the
      order of their literals are not distinguished, and a clause is also
      identified with the set of its literals.
\item A clause is \emph{Horn} when \(k\le 1\) and \emph{definite Horn}
      when \(k=1\): it has at most one, respectively exactly one,
      unnegated literal. These are the Horn clauses of
      \cref{def:p02-tractable-classes} and the definite clauses of
      \cref{def:p02-w02-p02-019-horn-parts}.
\item A theory is in \emph{clausal form} when all its formulas are
      clauses. A clausal theory all of whose clauses are Horn (definite
      Horn) is a \emph{Horn} (\emph{definite Horn}) theory.
\item A formula is in \emph{conjunctive normal form} (CNF) when it is a
      conjunction of clauses, and in \emph{disjunctive normal form} (DNF)
      when it is a disjunction \(D_1\lor\cdots\lor D_n\) of conjunctions
      \(D_i\) of literals.
\item The source also refers to a theory \(T\) as the conjunction of its
      formulas. Every conjunction of all members of \(T\), in any order and
      bracketing, with \(\topv\) for empty \(T\), has the models
      \(\operatorname{Mod}_V(T)\), so this reading changes neither
      consistency nor consequence. Read in this way, a theory in clausal
      form is a CNF formula.
\item The source writes \(\phi(V)\) for a truth assignment to \(V\) and
      \(V_\phi(F)\) for the truth value of a formula \(F\) under it. These
      are the valuations \(v\in\mathsf{Val}(V)\) and formula values of
      \cref{def:p16-abduction-models}.
\item If \(E\) is a formula and \(\phi(X)\) is a truth assignment to a set
      \(X=\{x_1,\ldots,x_n\}\) of variables occurring in \(E\), then
      \(E_{\phi(X)}\) is the formula obtained from \(E\) by replacing each
      occurrence of each \(x\in X\) by \(\topv\) when \(\phi(x)\) is true
      and by \(\bottom\) when \(\phi(x)\) is false. In this item only,
      \(E\) is the source's letter for a formula; elsewhere in the chapter
      \(E\) is a candidate.
\end{enumerate}
\end{hbdefinition}

\begin{hbtheorem}{thm:p15-w02c-p15-027-theory-forms}{proved=annex:p15-w02c-p15-027-theory-forms}{Theory forms and the explanation families}
Let \(\mathcal P=(V,H,M,T)\) be a finite abduction input.
\begin{enumerate}
\item If \(T'\) is a finite set of formulas over \(V\) with
      \(\operatorname{Mod}_V(T')=\operatorname{Mod}_V(T)\), then
      \(\mathcal P'=(V,H,M,T')\) has the same explanations,
      subset-minimal explanations and minimum-cardinality explanations
      as \(\mathcal P\).
\item For \(w\in\mathsf{Val}(V)\) let \(C_w\) be the clause
      \(\{\neg p:p\in V,\ w(p)=\mathbf T\}\cup\{p:p\in V,\ w(p)=\mathbf F\}\).
      Let \(W\) be the set of \(w\in\mathsf{Val}(V)\) that falsify some
      member of \(T\). The clausal theory
      \(T^{\rm cl}=\{C_w:w\in W\}\) is over \(V\), and
      \[
       \operatorname{Mod}_V(T^{\rm cl})=\operatorname{Mod}_V(T).
      \]
      Hence every input has a clausal-form version with the same
      explanation families.
\item Let \(E\) be a formula over \(V\), \(X\) a set of variables occurring
      in \(E\), and \(v\in\mathsf{Val}(V)\) with \(\phi(X)\) the restriction
      of \(v\) to \(X\). Then \(E\) and \(E_{\phi(X)}\) have the same value
      under \(v\).
\item The theories of the six inputs in \cref{def:p16-abduction-examples}
      that contain implications have the following clausal equivalents
      with the same models over the same alphabets:
      \[
      \begin{aligned}
       \mathcal P_{\rm sizes}&:\ \{m\lor\neg a,\ m\lor\neg b\lor\neg c\},\\
       \mathcal P_{\rm constraint}&:\ \{m\lor\neg a,\ \neg a\lor\neg b\},\\
       \mathcal P_{\rm permitted}&:\ \{m\lor\neg z\}.
      \end{aligned}
      \]
      The first and third are definite Horn theories. The second is a Horn
      theory that is not definite Horn, since \(\neg a\lor\neg b\) has
      \(k=0\).
\end{enumerate}
\end{hbtheorem}

For example, with \(E=(a\to m)\) and \(\phi(\{a\})\) assigning true to
\(a\), the substitution gives \(E_{\phi(\{a\})}=(\topv\to m)\), which has
the value of \(m\) under every valuation that makes \(a\) true. The
clausal theory of item~2 can be much longer than \(T\); the theorem
concerns the explanation families only and makes no claim about size.

\begin{hbtheorem}{thm:p15-w02c-p15-027-formula-hypotheses}{proved=annex:p15-w02c-p15-027-formula-hypotheses}{Formula hypotheses reduce to atoms}
Let \(V\) be finite, \(M\subseteq V\), \(T\) a finite set of formulas over
\(V\), and \(\Phi\) a finite set of formulas over \(V\), the permitted
\emph{formula hypotheses}. Call \(S\subseteq\Phi\) a \emph{formula
explanation} when \(\operatorname{Mod}_V(T\cup S)\ne\varnothing\) and
\(T\cup S\entails M\). Choose distinct atoms \(x_\varphi\notin V\) of
\(\mathsf{Prop}\), one for each \(\varphi\in\Phi\), and put
\[
\begin{aligned}
 V'&=V\cup\{x_\varphi:\varphi\in\Phi\},\qquad
 H'=\{x_\varphi:\varphi\in\Phi\},\\
 T'&=T\cup\{x_\varphi\leftrightarrow\varphi:\varphi\in\Phi\},\qquad
 \mathcal P'=(V',H',M,T').
\end{aligned}
\]
Then:
\begin{enumerate}
\item \(T'\) is consistent exactly when \(T\) is, so \(\mathcal P'\)
      meets the source's consistency requirement whenever \(T\) does.
\item The map \(S\mapsto\{x_\varphi:\varphi\in S\}\) is a bijection
      from the formula explanations onto \(\operatorname{Expl}(\mathcal P')\)
      that preserves inclusion and cardinality. Hence it carries the
      subset-minimal formula explanations onto
      \(\operatorname{MinExpl}(\mathcal P')\) and those of minimum
      cardinality onto \(\operatorname{CardExpl}(\mathcal P')\).
\end{enumerate}
Negative literals are the case in which a member of \(\Phi\) is
\(\neg p\) for an atom \(p\).
\end{hbtheorem}

The translation preserves the explanation relation, but it changes
what is counted. Two logically equivalent members of \(\Phi\) receive different
new atoms and so count as different hypotheses, and minimality is judged
over the chosen set \(\Phi\), not over all formulas.

Eiter and Gottlob's own construction is different. It keeps the
hypothesis atoms of \(H\) and adds one new atom \(x'\) for a formula
\(F\), with \(x'\leftrightarrow F\) added to the theory
\cite[manuscript p.~10]{EiterGottlobAbductionManuscript}. The next theorem
gives its exact effect, for one formula as in the source and for finitely
many added at once.

\begin{hbtheorem}{thm:p15-w02c-p15-027-source-extension}{proved=annex:p15-w02c-p15-027-source-extension}{Adding formula hypotheses beside the atoms of \texorpdfstring{\(H\)}{H}}
Let \(\mathcal P=(V,H,M,T)\) be a finite abduction input and \(\Phi\) a
finite set of formulas over \(V\). Choose distinct atoms
\(x'_\varphi\notin V\) of \(\mathsf{Prop}\), one for each
\(\varphi\in\Phi\); for \(\Psi\subseteq\Phi\) put
\(X_\Psi=\{x'_\varphi:\varphi\in\Psi\}\), and let
\[
 \mathcal P^{\Phi}=\bigl(V\cup X_\Phi,\ H\cup X_\Phi,\ M,\
   T^{\Phi}\bigr),\qquad
 T^{\Phi}=T\cup\{x'_\varphi\leftrightarrow\varphi:\varphi\in\Phi\}.
\]
Call a pair \((S,\Psi)\) with \(S\subseteq H\) and \(\Psi\subseteq\Phi\) a
\emph{mixed explanation} when
\(\operatorname{Mod}_V(T\cup S\cup\Psi)\ne\varnothing\) and
\(T\cup S\cup\Psi\entails M\).
\begin{enumerate}
\item \(T^{\Phi}\) is consistent exactly when \(T\) is.
\item The map \((S,\Psi)\mapsto S\cup X_\Psi\) is a bijection from the
      mixed explanations onto \(\operatorname{Expl}(\mathcal P^{\Phi})\).
      Moreover \(S_1\cup X_{\Psi_1}\subseteq S_2\cup X_{\Psi_2}\) exactly
      when \(S_1\subseteq S_2\) and \(\Psi_1\subseteq\Psi_2\), and
      \(|S\cup X_\Psi|=|S|+|\Psi|\).
\item The old explanations are kept, and so are the old subset-minimal
      ones:
      \(\operatorname{Expl}(\mathcal P)=\{E\in\operatorname{Expl}(\mathcal P^{\Phi}):E\subseteq H\}\)
      and
      \(\operatorname{MinExpl}(\mathcal P)=\{E\in\operatorname{MinExpl}(\mathcal P^{\Phi}):E\subseteq H\}\).
\item In the source's case \(\Phi=\{F\}\), with \(x'=x'_F\), let
      \(\mathcal S_F\) be the family of sets \(S\subseteq H\) with
      \(\operatorname{Mod}_V(T\cup S\cup\{F\})\ne\varnothing\) and
      \(T\cup S\cup\{F\}\entails M\). Then
      \[
       \operatorname{Expl}(\mathcal P^{\{F\}})
        =\operatorname{Expl}(\mathcal P)\cup
         \{S\cup\{x'\}:S\in\mathcal S_F\},
      \]
      and \(S\cup\{x'\}\) is subset-minimal exactly when \(S\) is
      inclusion-minimal in \(\mathcal S_F\) and no subset of \(S\)
      belongs to \(\operatorname{Expl}(\mathcal P)\).
\item Minimum cardinality is not preserved. Let
      \[
       \mathcal P=(\{a,b,m\},\{a,b\},\{m\},\{(a\land b)\to m\}),
       \qquad F=a\land b.
      \]
      Then
      \(\operatorname{Expl}(\mathcal P)=\operatorname{MinExpl}(\mathcal P)
        =\operatorname{CardExpl}(\mathcal P)=\{\{a,b\}\}\), while
      \[
      \begin{aligned}
       \operatorname{Expl}(\mathcal P^{\{F\}})
        &=\{\{a,b\},\{x'\},\{a,x'\},\{b,x'\},\{a,b,x'\}\},\\
       \operatorname{MinExpl}(\mathcal P^{\{F\}})&=\{\{a,b\},\{x'\}\},\qquad
       \operatorname{CardExpl}(\mathcal P^{\{F\}})=\{\{x'\}\}.
      \end{aligned}
      \]
\end{enumerate}
\end{hbtheorem}

So the source's claim that the atom convention loses no generality holds
for the explanation relation, in the form of items~2 and~4: every
explanation that uses the formula \(F\) is represented by exactly one
explanation of the extended input, and the old explanations remain. It
does not hold for every preference. In the example of item~5, \(a\) and
\(b\) are \(\subseteq\)-necessary for \(\mathcal P\) in the sense of
\cref{def:p15-w02c-p15-027-relevance}, but only \(\subseteq\)-relevant
for \(\mathcal P^{\{F\}}\); they are \(\le\)-necessary for
\(\mathcal P\) and \(\le\)-irrelevant for \(\mathcal P^{\{F\}}\). A question about relevance or necessity is therefore posed on
a different input once formula hypotheses are added. When \(H\) is
empty, \(\mathcal P^{\Phi}\) is the input \(\mathcal P'\) of
\cref{thm:p15-w02c-p15-027-formula-hypotheses} after renaming the new
atoms.

\begin{hbwarning}{def:p16-abduction-two-tests}{Consistency and sufficiency are separate tests}
\par\nobreak
Classical entailment from an inconsistent set is vacuous: there is no
satisfying valuation to falsify an observation. Such a set fails the
consistency test and gives no explanation. Conversely, consistency only
provides one model of the proposed assumptions. Explanation requires every
such model to satisfy every observation. A single countervaluation refutes
that requirement.

Membership in \(H\) is an additional input restriction. A consistent set
of atoms that entails \(M\) can still be ineligible because it includes
an atom outside \(H\). The examples below distinguish these failures.
\end{hbwarning}

\section{Two notions of a small explanation}
\label{sec:p16-abduction-minimality}

\begin{hbdefinition}{def:p16-abduction-minimality}{Subset-minimal and minimum-cardinality explanations}
An explanation \(E\) is \emph{subset-minimal} if no proper subset of
\(E\) is an explanation. It has \emph{minimum cardinality} if no
explanation has fewer members. Define the two families by
\[
\begin{aligned}
 \operatorname{MinExpl}(\mathcal P)
   &\defeq\{E\in\operatorname{Expl}(\mathcal P):\\
   &\hspace{1.5em}\text{no }F\subsetneq E\text{ is an explanation}\},\\
 \operatorname{CardExpl}(\mathcal P)
   &\defeq\{E\in\operatorname{Expl}(\mathcal P):\\
   &\hspace{1.5em}|E|\leq|F|\text{ for every explanation }F\}.
\end{aligned}
\]
If there are no explanations, both families are empty. These are the
subset and cardinality preference criteria of the source
\cite[def.~2.3 and following discussion, manuscript pp.~10--11]{EiterGottlobAbductionManuscript};
\cref{thm:p15-w02c-p15-027-preference} derives both families from the
source's general definition of preferred solutions.
They compare sets of assumptions. No probability, cost or causal relation
has been specified by either definition.
\index{abduction!subset-minimal explanation}
\index{abduction!minimum-cardinality explanation}
\end{hbdefinition}

\begin{hbtheorem}{thm:p16-abduction-minimality}{proved=annex:p16.abduction-minimality}{Existence, inclusion and the deletion criterion}
\par\nobreak
Every explanation contains a subset-minimal explanation. If any explanation
exists, both minimality families are nonempty, and
\[
 \operatorname{CardExpl}(\mathcal P)
        \subseteq\operatorname{MinExpl}(\mathcal P)
        \subseteq\operatorname{Expl}(\mathcal P).
\]
Distinct subset-minimal explanations are incomparable under inclusion;
their cardinalities may differ.

If \(E\) is an explanation, then
\[
\begin{aligned}
 E\in\operatorname{MinExpl}(\mathcal P)
 \quad\Longleftrightarrow\quad&
 E\setminus\{h\}\notin\operatorname{Expl}(\mathcal P)
                  \\
 &\text{for every }h\in E.
\end{aligned}
\]
This is the source's single-deletion characterization
\cite[prop.~2.1, p.~11]{EiterGottlobAbductionManuscript}.
It also holds for the extended input convention used here.

The relevant monotonicity has a consistency qualification: if \(E\) is
an explanation, \(E\subseteq F\subseteq H\), and \(T\cup F\) is
consistent, then \(F\) is an explanation. An arbitrary superset of an
explanation need not be an explanation, because it can become inconsistent.
\end{hbtheorem}

Subset-minimality can therefore be checked by testing each one-atom
deletion of a known explanation. Minimum cardinality compares it with all
explanations, including those incomparable with it. The unequal-size example
below makes the distinction explicit.

\section{Preference orders, irredundancy and the main problems}
\label{sec:p15-w02c-p15-027-preference-orders}

Eiter and Gottlob obtain the two families of
\cref{def:p16-abduction-minimality} as instances of one construction:
a preference relation on sets of hypotheses selects the acceptable
solutions. They take reflexivity and transitivity to be the natural
requirements on such a relation, and therefore model preferences by
preorders \cite[manuscript p.~10]{EiterGottlobAbductionManuscript}.
In this section \(\operatorname{Sol}(\mathcal P)\) is the source's name
for \(\operatorname{Expl}(\mathcal P)\).

\begin{hbdefinition}{def:p15-w02c-p15-027-preference}{Preference preorders, preferred solutions and irredundancy}
Let \(\mathcal P=(V,H,M,T)\) be a finite abduction input
\cite[defs.~2.3--2.4, manuscript pp.~10--11]{EiterGottlobAbductionManuscript}.
\begin{enumerate}
\item A \emph{preference} is a preorder \(\preceq\) on the powerset
      \(2^H\), that is, a reflexive and transitive relation. Its
      \emph{strict part} is defined by
      \(S\prec S'\) if and only if \(S\preceq S'\) and \(S'\not\preceq S\).
\item The \emph{preferred}, or \emph{acceptable}, solutions under
      \(\preceq\) are the \(\prec\)-minimal solutions:
      \[
       \operatorname{Sol}_{\preceq}(\mathcal P)
        \defeq\{S\in\operatorname{Sol}(\mathcal P):
          \text{there is no }S'\in\operatorname{Sol}(\mathcal P)
          \text{ with }S'\prec S\}.
      \]
\item Three preferences are used: equality \(=\); the
      \emph{subset-minimality order} \(\subseteq\), with \(S\preceq S'\)
      if and only if \(S\subseteq S'\); and the \emph{minimum solution size}
      (minimum cardinality) order \(\le\), with \(S_1\le S_2\) if and only
      if \(|S_1|\le|S_2|\). For the size order the source cites Peng and
      Reggia's book on abductive diagnosis (its reference~[59]).
\item The preference \(\preceq\) is \emph{irredundant} if and only if,
      for all \(S,S'\in\operatorname{Sol}(\mathcal P)\), \(S\subset S'\)
      implies \(S\prec S'\). Here \(\subset\) is proper inclusion; with
      \(S=S'\) the condition would demand \(S\prec S\), which no preorder
      satisfies. The source attributes the importance of this property to
      Peng and Reggia, Selman and Levesque, Bylander, Allemang, Tanner and
      Josephson, and Konolige (its references~[59], [69], [6], [40]).
      The members of \(\operatorname{Sol}_{\subseteq}(\mathcal P)\) are
      called the \emph{irredundant solutions}; the source describes the
      subset-minimality order as irredundancy itself.
\end{enumerate}
The definition of irredundancy quantifies over the solutions of the given
input, so it is a property of a preference relative to \(\mathcal P\).
\index{abduction!preference preorder}\index{abduction!irredundant solution}
\end{hbdefinition}

\begin{hbtheorem}{thm:p15-w02c-p15-027-preference}{proved=annex:p15-w02c-p15-027-preference}{Strict parts and irredundancy}
Let \(\mathcal P=(V,H,M,T)\) be a finite abduction input.
\begin{enumerate}
\item The relations \(=\), \(\subseteq\) and \(\le\) are preorders on
      \(2^H\). Their strict parts are respectively the empty relation,
      proper inclusion \(S\subsetneq S'\), and \(|S|<|S'|\). Consequently
      \[
      \begin{aligned}
       \operatorname{Sol}_{=}(\mathcal P)&=\operatorname{Expl}(\mathcal P),\\
       \operatorname{Sol}_{\subseteq}(\mathcal P)&=\operatorname{MinExpl}(\mathcal P),\\
       \operatorname{Sol}_{\le}(\mathcal P)&=\operatorname{CardExpl}(\mathcal P).
      \end{aligned}
      \]
      In particular \(\operatorname{Sol}(\mathcal P)=\operatorname{Sol}_{=}(\mathcal P)\).
\item For every preorder \(\preceq\) on \(2^H\), the strict part \(\prec\)
      is irreflexive and transitive, and
      \(\operatorname{Sol}_{\preceq}(\mathcal P)\) is nonempty exactly when
      \(\operatorname{Expl}(\mathcal P)\) is nonempty.
\item The orders \(\subseteq\) and \(\le\) are irredundant for every input.
      Equality is irredundant for \(\mathcal P\) exactly when no solution
      is a proper subset of another; it is not irredundant for
      \(\mathcal P_{\rm given}\), whose solutions are
      \(\varnothing\subsetneq\{a\}\).
\item If \(\preceq\) is irredundant for \(\mathcal P\), then
      \(\operatorname{Sol}_{\preceq}(\mathcal P)
        \subseteq\operatorname{Sol}_{\subseteq}(\mathcal P)\).
\end{enumerate}
The irredundant solutions are thus the subset-minimal explanations, and
\cref{thm:p16-abduction-minimality} gives their single-deletion
characterization.
\end{hbtheorem}

Item~4 explains the role of irredundancy: an irredundant preference never
accepts a solution that has a solution as a proper subset. The size order
can be strictly more selective than subset-minimality, as
\(\mathcal P_{\rm sizes}\) shows: \(\{b,c\}\) is irredundant but not of
minimum size.

\begin{hblemma}{thm:p15-w02c-p15-027-decomposition}{proved=annex:p15-w02c-p15-027-decomposition}{Inputs on disjoint alphabets decompose}
Let \(n\ge 1\) and let \(\mathcal P_i=(V_i,H_i,M_i,T_i)\), for
\(1\le i\le n\), be finite abduction inputs with
\(V_i\cap V_j=\varnothing\) for \(1\le i<j\le n\). Put
\(V=\bigcup_iV_i\), \(H=\bigcup_iH_i\), \(M=\bigcup_iM_i\),
\(T=\bigcup_iT_i\) and \(\mathcal P=(V,H,M,T)\). If \(\preceq\) is
equality, \(\subseteq\) or \(\le\), then
\[
 \operatorname{Sol}_{\preceq}(\mathcal P)
  =\{S_1\cup\cdots\cup S_n:
       S_i\in\operatorname{Sol}_{\preceq}(\mathcal P_i),\ 1\le i\le n\},
\]
and each member of the left side has exactly one such representation,
with \(S_i=S\cap H_i\). The disjointness hypothesis cannot be dropped:
for \(\mathcal P_1=(\{a,m\},\{a\},\{m\},\{a\to m\})\) and
\(\mathcal P_2=(\{a\},\varnothing,\varnothing,\{\neg a\})\), which share
the atom \(a\), the union input has no solution, while the right side is
\(\{\{a\}\}\).
\end{hblemma}

This is the source's Lemma~2.2, which it states for consistent theories
without a written proof
\cite[lem.~2.2, manuscript p.~11]{EiterGottlobAbductionManuscript}. The
proof in the annex also covers inconsistent \(T_i\).

\begin{hbdefinition}{def:p15-w02c-p15-027-relevance}{Existence, relevance and necessity}
In this definition \(\mathcal P=(V,H,M,T)\) is a finite abduction input,
\(\preceq\) a preference on \(2^H\) and \(h\in H\). The source's three main problems
are: whether \(\mathcal P\) has a solution (\emph{existence}); whether
\(h\) contributes to some acceptable solution (\emph{relevance}); and
whether \(h\) occurs in all acceptable solutions (\emph{necessity})
\cite[def.~2.5 and the preceding list, manuscript p.~11]{EiterGottlobAbductionManuscript}.
The source records that the relevance and necessity questions had been
studied before: by Bylander, Allemang, Tanner and Josephson (its
reference~5, Proceedings KR-89) for set-oriented models of abduction,
and, for particular subclasses of propositional abduction, by Selman and
Levesque (reference~69, Proceedings AAAI-90), by Friedrich, Gottlob and
Nejdl (reference~28, 1990) and by Bylander (reference~4, Proceedings
KR-91)
\cite[manuscript p.~11, and reference list, manuscript pp.~41, 42 and~45]{EiterGottlobAbductionManuscript}.
That is the source's historical report; the earlier works are not
examined here.
Precisely:
\begin{enumerate}
\item \(h\) is \emph{\(\preceq\)-relevant} for \(\mathcal P\) if and only
      if \(h\in S\) for some \(S\in\operatorname{Sol}_{\preceq}(\mathcal P)\);
\item \(h\) is \emph{\(\preceq\)-necessary} for \(\mathcal P\) if and only
      if \(h\in S\) for every \(S\in\operatorname{Sol}_{\preceq}(\mathcal P)\);
\item \(h\) is \emph{\(\preceq\)-irrelevant} if it is not
      \(\preceq\)-relevant, and \emph{\(\preceq\)-dispensable} if it is not
      \(\preceq\)-necessary. Josephson, Chandrasekaran, Smith and Tanner
      call necessity \emph{indispensability} (the source's reference~[36]).
\item When \(\preceq\) is equality, the prefix is dropped: \emph{relevant},
      \emph{necessary}, and so on.
\end{enumerate}
The source states that these notions divide \(H\) into three disjoint
sets: the hypotheses that are not \(\preceq\)-relevant, those that are
\(\preceq\)-relevant but not \(\preceq\)-necessary, and the
\(\preceq\)-necessary ones. It works mainly with dispensability, and
states that its results for necessity are easy corollaries of those for
dispensability. The source's subject is the computational complexity of
these decision problems; that classification is not developed here.
\index{abduction!relevance}\index{abduction!necessity}
\index{abduction!dispensability}
\end{hbdefinition}

\begin{hbtheorem}{thm:p15-w02c-p15-027-relevance}{proved=annex:p15-w02c-p15-027-relevance}{When relevance and necessity partition the hypotheses}
Let \(\mathcal P=(V,H,M,T)\) be a finite abduction input and \(\preceq\) a
preorder on \(2^H\). Let \(I\), \(R\) and \(N\) be the sets of hypotheses
that are \(\preceq\)-irrelevant, \(\preceq\)-relevant but not
\(\preceq\)-necessary, and \(\preceq\)-necessary.
\begin{enumerate}
\item For each \(h\in H\), \(h\) is \(\preceq\)-necessary exactly when it
      is not \(\preceq\)-dispensable. So, on each pair \((\mathcal P,h)\),
      the answer to the necessity question is the negation of the answer
      to the dispensability question.
\item Acceptable solutions exist exactly when solutions exist.
\item If \(\operatorname{Sol}_{\preceq}(\mathcal P)\ne\varnothing\), every
      \(\preceq\)-necessary hypothesis is \(\preceq\)-relevant, and \(I\),
      \(R\), \(N\) are pairwise disjoint with union \(H\).
\item If \(\operatorname{Sol}_{\preceq}(\mathcal P)=\varnothing\), then
      \(I=N=H\) and \(R=\varnothing\). Hence for \(H\ne\varnothing\) the
      three sets are not pairwise disjoint. This happens in
      \(\mathcal P_{\rm permitted}\), whose hypothesis \(a\) is both
      \(\preceq\)-irrelevant and \(\preceq\)-necessary for every
      preorder \(\preceq\).
\item In \(\mathcal P_{\rm sizes}\), under equality and under \(\subseteq\)
      each of \(a,b,c\) is relevant and none is necessary. Under \(\le\),
      \(a\) is \(\le\)-necessary and \(b,c\) are \(\le\)-irrelevant.
\end{enumerate}
\end{hbtheorem}

The source's division of \(H\) into three disjoint sets is therefore a
partition only when acceptable solutions exist. Item~5 shows that the
choice of preference changes relevance: \(b\) and \(c\) occur in the
subset-minimal explanation \(\{b,c\}\) and are relevant for \(\subseteq\),
but are irrelevant for the size order.

\section{An exhaustive finite construction}
\label{sec:p16-abduction-enumeration}

\begin{hbmethod}{def:p16-abduction-enumeration}{Enumerating explanations and both minimality families}
The input is the finite tuple \(\mathcal P=(V,H,M,T)\). List \(V\)
and \(H\) by increasing indices in \(\mathsf{Prop}=\{p_i:i\in\N\}\).
Put \(n=|V|\), \(r=|H|\), and initialize a list \(\mathcal L\)
with no members.
\begin{enumerate}
\item Enumerate every binary vector of length \(r\), in lexicographic
      order. Its candidate \(E\subseteq H\) contains the hypothesis
      at a position exactly when the bit at that position is one.
\item For this \(E\), initialize two Boolean flags to false:
      \emph{has a model} and \emph{has a countervaluation}. Enumerate
      every binary vector of length \(n\), in lexicographic order,
      interpreting zero as
      \(\mathbf{F}\) and one as \(\mathbf{T}\). This gives a valuation on
      \(V\). Evaluate every formula of \(T\) and every atom of \(E\).
      If all are true, set \emph{has a model} to true. If in addition some
      \(m\in M\) is false, set \emph{has a countervaluation} to true.
      Once set, a flag stays true for this candidate.
\item After the valuation scan, append \(E\) to \(\mathcal L\)
      exactly when \emph{has a model} is true and
      \emph{has a countervaluation} is false.
\item After all candidates, form \(\mathcal L_{\subseteq}\) by
      retaining exactly those members \(E\) of \(\mathcal L\) for
      which no other member is a proper subset of \(E\).
\item If \(\mathcal L\) is empty, put \(\mathcal L_{\#}\) empty.
      Otherwise compute the least cardinality of a member of \(\mathcal L\)
      by a finite scan, and let \(\mathcal L_{\#}\) retain exactly
      the members of that cardinality. Return all three lists.
\end{enumerate}
The empty binary vector occurs once when its length is zero. In particular,
empty \(H\) still gives the candidate \(\varnothing\), and empty
\(V\) still gives the unique empty valuation. Testing all observations
when \(M\) is empty cannot set the countervaluation flag.
\end{hbmethod}

\begin{hbtheorem}{thm:p16-abduction-enumeration}{proved=annex:p16.abduction-enumeration}{Termination and exact output}
The construction terminates. Viewed as sets of their listed members, its
outputs are exactly
\[
\begin{aligned}
 \mathcal L&=\operatorname{Expl}(\mathcal P),\\
 \mathcal L_{\subseteq}&=\operatorname{MinExpl}(\mathcal P),\\
 \mathcal L_{\#}&=\operatorname{CardExpl}(\mathcal P).
\end{aligned}
\]
No candidate or output member is duplicated. The unshortened nested scans
in steps 1--3 inspect exactly \(2^{r+n}\) candidate--valuation pairs.
This count does not make formula evaluation or the subsequent filtering
constant-time operations. No optimality or complexity-class claim is made.
\end{hbtheorem}

\section{Complete small examples}
\label{sec:p16-abduction-examples}

\begin{hbdefinition}{def:p16-abduction-examples}{Six explicit inputs}
In the following inputs, distinct letters among \(a,b,c,z,m\) denote
distinct members of \(\mathsf{Prop}\). The tuple order is always \((V,H,M,T)\).
These are additional handbook examples.
\[
\begin{aligned}
 \mathcal P_{\rm inconsistent}
   &\defeq(\{a,m\},\{a\},\{m\},\{a,\neg a\}),\\
 \mathcal P_{\rm given}
   &\defeq(\{a,m\},\{a\},\{m\},\{m\}),\\
 \mathcal P_{\rm empty}
   &\defeq(\varnothing,\varnothing,\varnothing,\varnothing),\\
 \mathcal P_{\rm sizes}
   &\defeq(\{a,b,c,m\},\{a,b,c\},\{m\},\\
   &\hspace{2em}\{a\to m,(b\land c)\to m\}),\\
 \mathcal P_{\rm constraint}
   &\defeq(\{a,b,m\},\{a,b\},\{m\},\\
   &\hspace{2em}\{a\to m,\neg(a\land b)\}),\\
 \mathcal P_{\rm permitted}
   &\defeq(\{a,z,m\},\{a\},\{m\},\{z\to m\}).
\end{aligned}
\]
\end{hbdefinition}

\begin{hbtheorem}{thm:p16-abduction-examples}{proved=annex:p16.abduction-examples}{All explanation families for the six inputs}
The complete families are as follows.
\begin{enumerate}
\item \emph{Inconsistent background.} All three explanation families of
      \(\mathcal P_{\rm inconsistent}\) are empty, even though every
      \(T\cup E\) entails \(m\) vacuously.
\item \emph{Observation already given.}
      \(\operatorname{Expl}(\mathcal P_{\rm given})
          =\{\varnothing,\{a\}\}\).
      Both minimality families equal \(\{\varnothing\}\).
\item \emph{Empty data.} All three families of
      \(\mathcal P_{\rm empty}\) equal \(\{\varnothing\}\).
      This is one explanation, rather than no explanations.
\item \emph{Unequal sizes.}
      \[
      \begin{aligned}
       \operatorname{Expl}(\mathcal P_{\rm sizes})
         &=\{\{a\},\{a,b\},\{a,c\},\{a,b,c\},\{b,c\}\},\\
       \operatorname{MinExpl}(\mathcal P_{\rm sizes})
         &=\{\{a\},\{b,c\}\},\\
       \operatorname{CardExpl}(\mathcal P_{\rm sizes})
         &=\{\{a\}\}.
      \end{aligned}
      \]
      The two subset-minimal explanations are incomparable, with sizes
      one and two.
\item \emph{Consistency constraint.} All three families of
      \(\mathcal P_{\rm constraint}\) equal \(\{\{a\}\}\).
      The candidate \(\{b\}\) is consistent with \(T\) but does
      not entail \(m\). The superset \(\{a,b\}\) of the explanation
      \(\{a\}\) fails consistency.
\item \emph{Permitted hypotheses.} All three families of
      \(\mathcal P_{\rm permitted}\) are empty. The set \(\{z\}\)
      is consistent with \(T\) and entails \(m\) with \(T\), but
      is not a candidate because \(z\notin H\). The permitted
      nonempty candidate \(\{a\}\) is consistent and non-explanatory.
\end{enumerate}
\end{hbtheorem}

\section{What an explanation establishes}
\label{sec:p16-abduction-interpretation}

\begin{hbtheorem}{thm:p16-abduction-truth-boundary}{proved=annex:p16.abduction-truth-boundary}{Explanatory assumptions need not follow from the observations}
In \(\mathcal P_{\rm sizes}\), both \(\{a\}\) and \(\{b,c\}\)
are subset-minimal explanations, but
\[
\begin{aligned}
 T\cup M&\not\entails a,\\
 T\cup M&\not\entails(b\land c).
\end{aligned}
\]
Indeed there is one valuation satisfying \(T\cup M\) which makes all
three permitted hypotheses false.
\end{hbtheorem}

An explanation asserts that, \emph{if} its hypotheses hold together with
the background, the observations follow, and that these combined assumptions
are satisfiable. The theorem shows why this does not establish that the
selected hypotheses are true. A declaration that \(M\) contains the
observations does not by itself supply evidence for their measurement or for
the background theory either.

\begin{hbwarning}{def:p16-abduction-causal-boundary}{Logical explanation and causal interpretation}
The arrows in \(T\) are classical material implications. The input
contains no causal relation, intervention semantics, temporal ordering or
probability distribution. Consequently the construction establishes a
logical property of the supplied formulas and permitted assumptions. A
claim about what actually caused an observation requires additional
structure and evidence. Choosing fewer assumptions, or removing redundant
ones, does not supply that structure.

\Cref{thm:p15-w02c-p15-027-causal-boundary} works a boundary case: two
causal models agree on every valuation the background allows and differ
on what an intervention on an explanation's hypotheses does.

This chapter supplies the finite atomic-hypothesis foundation and its stated
procedures; formula hypotheses reduce to it by
\cref{thm:p15-w02c-p15-027-formula-hypotheses}, and the source's
general preference framework is in
\cref{sec:p15-w02c-p15-027-preference-orders}. First-order
abduction, preference orders other than the three treated there, the
source's complexity classifications,
probabilistic abduction and causal explanation require separate treatments.
Consistency-based diagnosis, its complete conflict search, measurement
updates and default correspondence are supplied in
\cref{sec:p16-consistency-based-diagnosis}.
The exact distinction between consistency diagnosis and explanatory
entailment is proved in \cref{thm:p16-diag-abduction-boundaries}.
The distinction between generating an analogical candidate and
establishing its truth by entailment is
\cref{thm:p19-analogy-candidate-truth}; the associated task evidence
is supplied in \cref{sec:p19-analogy-and-problem-transfer}.
The separate deductive-explanation model, its verification
condition and its account of relative emergence are supplied in
\cref{sec:p16-deductive-nomological-explanation}.
The complete source
presentation and the broader explanation/abduction subject remain beyond
this bounded unit.
\end{hbwarning}

\begin{hbtheorem}{thm:p15-w02c-p15-027-causal-boundary}{proved=annex:p15-w02c-p15-027-causal-boundary}{The same explanations under opposite causal readings}
Identify the atoms \(a,b,c,m\) of \(\mathcal P_{\rm sizes}\) with the
endogenous vertices of two finite structural causal models in the sense
of \cref{def:p16-finite-scm}, each with value set \(\mathbb{B}\).
In both, the outcome space is \(\Omega=\mathbb{B}^4\) with mass
\(1/16\) at each point, and the disturbances \(U_a,U_b,U_c,U_m\) are
the four coordinates. Model \(C_1\) has edges from \(a,b,c\) to \(m\)
and the equations
\[
 a=U_a,\quad b=U_b,\quad c=U_c,\quad m=a\lor(b\land c)\lor U_m.
\]
Model \(C_2\) has edges from \(m\) to \(a\) and \(b\) and from \(c\)
to \(b\), and the equations
\[
 m=U_m,\quad c=U_c,\quad a=m\land U_a,\quad b=U_b\land(m\lor\neg c).
\]
Then:
\begin{enumerate}
\item In each model the set of solutions attained at the outcomes is
      exactly \(\operatorname{Mod}_V(T)\) for the background \(T\) of
      \(\mathcal P_{\rm sizes}\), and each of its eleven valuations has
      positive probability. So both models are compatible with exactly the
      valuations the input allows, and its explanation families, given in
      \cref{thm:p16-abduction-examples}, are the same whichever model is
      intended.
\item In \(C_1\), the intervention \(\operatorname{do}(a=\mathbf T)\)
      and the intervention \(\operatorname{do}(b=\mathbf T,c=\mathbf T)\)
      each make \(m\) true at every outcome.
\item In \(C_2\), under either intervention \(m\) keeps its value
      \(U_m\) at every outcome, so \(m\) is false with probability
      \(1/2\).
\item The countervaluation of \cref{thm:p16-abduction-truth-boundary},
      with \(a,b,c\) false and \(m\) true, is a solution of positive
      probability in both models.
\end{enumerate}
\end{hbtheorem}

In \(C_1\), each subset-minimal explanation names conditions whose
production brings about the observation. In \(C_2\), the observation is
upstream of \(a\) and \(b\), and producing the hypotheses leaves it
untouched. The input, its explanations and their minimality are identical
in the two cases. Item~4 adds that in both models the observation can hold
with every hypothesis false. It happens at the outcome where \(U_m\) is
true and \(U_a,U_b,U_c\) are false. In \(C_1\), the hypotheses then take
the false values of their disturbances, and \(m\) is true through \(U_m\)
alone. In \(C_2\), \(m\) has no endogenous parent and its equation is
\(m=U_m\), so \(m\) is true; \(c=U_c\) is false; and \(a=m\land U_a\) and
\(b=U_b\land(m\lor\neg c)\) are false because \(U_a\) and \(U_b\) are
false. In both models \(m\) is an endogenous vertex; in the sense of
\cref{def:p16-finite-scm} only the disturbances are exogenous. The causal reading of the arrows in
\(T\) is therefore fixed by structure the abduction input does not
contain, as \cref{def:p16-abduction-causal-boundary} states.
